\section{Syntactic categories}
\label{sec:smt-syntax}

\subsection{SMT types}
\label{subsec:smt-type-system}

SMT-LIB is intrinsically many-sorted: every well-formed term has a sort, and the logic is defined
over a fixed signature of sorts and sorted function symbols.
We therefore introduce a syntactic category of SMT types, with two base sorts,
integers and Booleans, a unit sort, and three type constructors: functions,
options, and pairs.
Note that this type system is tailored to our specific formalization needs and does not
attempt to capture the full generality of SMT-LIB's sort system, which includes user-defined
sorts and polymorphic type constructors.

In order to simplify the type system even further, we make use of SMT-LIB's recent support
(version 2.7) for polymorphic data-types to encode pairs and options.
In the formalization, they are treated as primitive type constructors, and SMT encodings always
import a prelude containing those definitions. In practice, they are implemented as follows:

\noindent
\begin{minipage}[t]{0.45\textwidth}
  \begin{smtlisting}
(declare-datatype Pair
  (par (T1 T2)
    ((pair (fst T1) (snd T2)))))
  \end{smtlisting}
\end{minipage}\hfill
\begin{minipage}[t]{0.45\textwidth}
  \begin{smtlisting}
(declare-datatype Option
  (par (T)
    ((some (the T)) (none))))
  \end{smtlisting}
\end{minipage}
\medskip

In the following, we write \(\tau \timesS \sigma\) for the pair sort \smtinl{Pair $\tau$ $\sigma$} and \(\optS\tau\) for the option sort \smtinl{Option $\tau$}.
Note that since the term constructor \smtinl{none} is polymorphic, it must be annotated with its sort when used in practice, to allow solvers to determine its sort; this is achieved by using the operator \smtinl{as}, as illustrated in the following example.

\begin{example}[SMT-LIB pairs and options]
  \label{ex:smt-pairs-options}
  The following SMT-LIB script declares a pair of an integer and a Boolean \smtinl{p}, and an option of integers \smtinl{o}, and asserts that the first component of the pair is zero and the option is empty:
\begin{smtlisting}
(declare-const p (Pair Int Bool))
(declare-const o (Option Int))
(assert (= (fst p) 0))
(assert (= o (as none (Option Int))))
\end{smtlisting}
\end{example}

We then define the syntactic category of SMT types as the following inductive type.
\begin{definition}[SMT types][SMT.SMTType][SMT/Syntax.html\#SMT.SMTType]
  \label{def:smt-types}
  \begin{bnf}[colspec = {l@{}l@{\kern4pt}c@{\kern4pt}ll}]
    \leaninl{SMT.SMTType} \(\,\ni\, \tau, \sigma\) ::= \(\regnotation{\boolS}\) && \(\regnotation{\intS}\) && \(\regnotation{\unitS}\) && \(\regnotation{\tau \toS \sigma}\) && \(\regnotation{\optS}\tau\) && \(\tau \regnotation{\timesS}\sigma\)
  \end{bnf}
\end{definition}

Functions in SMT-LIB are first-class and total, so the function sort \(\tau \toS \sigma\) is the sort of total functions from \(\tau\) to \(\sigma\).
This is a crucial point of divergence from the source language, and options will be used to
model partiality.
The unit sort \(\unitS\) is a singleton sort, whose sole inhabitant is written \(\aop{\unitvalS}\).
It appears only in the abstract syntax, since the encoder never emits a bare unit literal in the
concrete syntax.

\subsection{Concrete terms}
\label{subsec:smt-concrete-terms}

We then define a syntactic category for concrete SMT terms, giving a
first-order representation of the fragment of SMT-LIB targeted by the encoding.
This deep embedding constitutes the concrete layer of the formalization in Lean.
Concrete variables are names of type \leaninl{String}.
\begin{definition}[Concrete terms][SMT.Term][SMT/Syntax.html\#SMT.Term]
  \label{def:smt-basic-syntax}
                                                The constructors of \leaninl{SMT.Term} are listed below, where \(\var{v}\) is a name
  in \leaninl{String}, \(n\) is an integer, \(b\) is a Boolean value, \(\tau\) is
  an SMT type, \(\vec{\var{v}}\) is a list of names, \(\vec\tau\) is a list of SMT
  types, and \(\vec t\) is a list of terms.
  As in \Cref{ch:b}, object variables are typeset in bold gray throughout.

  \begin{bnf}[colspec = {l@{}l@{\kern6pt}c@{\kern6pt}ll}]
    \leaninl{SMT.Term} $\,∋\, t, s$ ::=
    | $\var{v}$ && $\mathsf{int}\ n$ && $\mathsf{bool}\ b$ // Atoms
    | $\regnotation{\appS{t}{s}}$ // Higher-order application
    | $\regnotation{\lamS \vec{\var{v}} : \vec\tau \cdot t}$ && $\regnotation{\forallS \vec{\var{v}} : \vec\tau \cdot t}$ && $\regnotation{\existsS \vec{\var{v}} : \vec\tau \cdot t}$ // Binders
    | $t \regnotation{\eqS} s$ && $t \regnotation{\andS} s$ && $t \regnotation{\orS} s$  // Boolean rules
    | $\regnotation{\notS t}$ && $t \regnotation{\impS} s$ && $\regnotation{\iteS}\ t\ s\ r$
    | $\regnotation{\someS} t$ && $\regnotation{\theS} t$ && $\regnotation{\noneS}$ && $\regnotation{\asS}\ t\ \tau$ // Options
    | $\regnotation{\pairS{t}{s}}$ && $\regnotation{\fstS} t$ && $\regnotation{\sndS} t$ // Pairs
    | $\regnotation{\distinctS} \vec t$ // Distinctness
    | $t \regnotation{\leS} s$ && $t \regnotation{\addS} s$ && $t \regnotation{\subS} s$ && $t \regnotation{\mulS} s$ // Arithmetic
  \end{bnf}
\end{definition}

\subsection{Abstract syntax}
\label{subsec:smt-abstract-syntax}

The abstract syntax is given by the inductive family
{\small\leancodeptr{SMT.PHOAS.Term}{SMT/PHOAS/Def.html\#SMT.PHOAS.Term}}
parameterized by a type \(\mathcal{V}\) of abstract variables.
In the grammar below, \(\tau\) ranges over SMT types, \(\vec\tau\) over finite
families of SMT types, and \(e\) over scopes of the displayed type; the arity of a
binder is inferred from \(\vec\tau\).
The non-binding constructors are structurally the same as in \Cref{def:smt-basic-syntax}.
As in the source language, binding constructors represent the scope of bound
variables by meta-level Lean functions; we reuse the polyadic encoding of
\Cref{subsec:b-abstract-syntax}, in which a binder of arity \(n\) takes a type
vector and a scope
\(
  e : (\mathsf{Fin}\,n \to \mathcal{V}) \to \mathsf{Term}\,\mathcal{V}.
\)

\begin{definition}[Abstract syntax][SMT.PHOAS.Term][SMT/PHOAS/Def.html\#SMT.PHOAS.Term]
  \label{def:smt-abstract-syntax}
      The constructors for abstract terms parametrized by a type \(\mathcal{V}\) of abstract variables
  are listed below, keeping notations from \Cref{def:smt-basic-syntax}.

  \begin{bnf}[colspec = {l@{}l@{\kern6pt}c@{\kern6pt}ll}]
    \leaninl{SMT.PHOAS.Term} $\,∋\, t, s$ ::=
    | $\var{v}$ && $\mathsf{int}\ n$ && $\mathsf{bool}\ b$ // Atoms
    | $\appSa{t}{s}$ // Application
    | $\aop{\lamS} \vec\tau \cdot e$ && $\aop{\forallS} \vec\tau \cdot e$ // Binders
    | $t \abin{\eqS} s$ && $t \abin{\andS} s$ && $\aop{\notS} t$ && $\aop{\iteS}\ t\ s\ r$ // Boolean rules
    | $\someSa t$ && $\theSa t$ && $\aop{\noneS}_{\tau}$ && $\regnotation{\aop{\unitvalS}}$ // Options and unit
    | $\pairSa{t}{s}$ && $\fstSa t$ && $\sndSa t$ // Pairs
    | $\distinctSa \vec t$ // Distinctness
    | $t \abin{\leS} s$ && $t \abin{\addS} s$ && $t \abin{\subS} s$ && $t \abin{\mulS} s$ // Arithmetic
  \end{bnf}
\end{definition}

\begin{remark}
  As in \Cref{subsec:b-abstract-syntax}, the arity \(n\) of a binder is implicit
  and inferred from the type of its type family \(\vec \tau\) and scope \(e\),
  so binders need not be annotated with their arity when writing terms.
\end{remark}

Unlike the source language, the concrete syntax given in \Cref{def:smt-basic-syntax} is not minimal:
several logical connectives and the existential quantifier are definable from a smaller core,
and are treated as derived forms in the abstract syntax of \Cref{def:smt-abstract-syntax}.
The reason for this choice is that we want fine-grained control over the concrete layer, since it
is the one emitted by the encoder; the abstract layer, on the other hand, is only used at proof
level, and it is therefore preferable to keep it as small as possible to minimize the number
of cases in proofs.
Of course, this means that we have to factor out some of the concrete constructors into derived
forms, thereby involving meta-level (semantic) considerations upon which one might want to get
a better grip.
We therefore record the defining equations here, as they also fix the intended reading of the
corresponding concrete constructors.

\begin{definition}[Derived connectives][SMT.PHOAS.Term.or, .imp, .exists][
    SMT/PHOAS/Def.html\#SMT.PHOAS.Term.or,
    SMT/PHOAS/Def.html\#SMT.PHOAS.Term.imp,
    SMT/PHOAS/Def.html\#SMT.PHOAS.Term.exists]
  \label{def:smt-derived}
  We define the following derived connectives and quantifier:
  \begin{align*}
    t \abin{\orS} s                    & \triangleq \aop{\notS}((\aop{\notS} t) \abin{\andS} (\aop{\notS} s)) \tag{disjunction}      \\
    t \abin{\impS} s                   & \triangleq \aop{\notS}(t \abin{\andS} (\aop{\notS} s)) \tag{implication}                     \\
    \existsS \vec\tau \cdot P & \triangleq \aop{\notS}(\aop{\forallS} \vec{\var{v}} : \vec\tau \cdot \aop{\notS} P) \tag{existential quantification}
  \end{align*}
\end{definition}

We also get rid of the annotation construct \(\asS\) in the abstract syntax, since in our case, it is only used when resolving a \(\noneS\) type, which we make primitive as a constructor \(\aop{\noneS}_{\tau}\), in the following informal sense:
\[
  \aop{\noneS}_{\tau} \;\approx\; \asS\ \noneS\ (\optS\tau).
\]

\begin{note}
  In contrast with the bounded binders of B, which range over a domain
  set \(D\) (written \(\vec{\var{v}} \inB D\)), SMT binders are sorted:
  each bound variable \(\var{v_i}\) ranges over an entire sort \(\tau_i\) (written
  \(\vec{\var{v}} : \vec\tau\)).
  This reflects the syntactic counterpart of another semantic shift, from membership constraints to
  sort membership.
\end{note}

Substitution is defined directly on abstract terms; since the PHOAS
representation delegates \(\alpha\)-conversion and capture-avoidance to Lean's
meta-level binding infrastructure, the definition is homomorphic everywhere
except at variables and binders.

\begin{definition}[Substitution][SMT.PHOAS.subst][SMT/PHOAS/Def.html\#SMT.PHOAS.subst]
  \label{def:smt-subst}
  We write \(t[\var{v} \leftarrow s]\) for the substitution of the term \(s\) for the
  variable \(\var{v}\) in \(t\). For a variable \(\var{w} : \mathcal{V}\),
  \[
    (\var{w})[\var{v} \leftarrow s] \triangleq
    \begin{cases}
      s               & \text{if } \var{w} = \var{v}, \\
      \var{w} & \text{otherwise;}
    \end{cases}
  \]
  for a binder with type vector \(\vec\tau\) and scope \(e\),
  \[
    (\lamS \vec\tau \cdot e)[\var{v} \leftarrow s] \triangleq \lamS \vec\tau \cdot (\pink{\lambda} \vec x \mapsto (e\,\vec x)[\var{v} \leftarrow s]),
  \]
  and analogously for the universal binder; all remaining clauses are
  homomorphic on subterms.
  Note that the \(\pink{\lambda}\) is a Lean meta-level binding.
\end{definition}

\subsection{Changing representations}
\label{subsec:smt-syntax-change-repr}

As on the source side, concrete terms are interpreted only after abstraction to
PHOAS. The relationship between the two representations is the same as in
\Cref{fig:syntax-representation}: abstraction maps concrete variable names to
PHOAS variables, while reification recovers names from the meta-level scope.
Only abstraction is implemented, since reification is not needed in the encoding
pipeline.

Since definitions are essentially the same as on the source side, we do not unfold them here.
Abstraction of a concrete term \(t\) is defined structurally and also denoted \(\aterm{t}{\Delta}\),
where \(\Delta\) is a valuation defined at least on the free variables of \(t\).
\begin{note}
  In order to avoid any ambiguity, this notation will be used either under the denotation function
  \(\interpB{\cdot}\) or \(\interpS{\cdot}\) to explicitly indicate the underlying language, or on
  a term that will be properly contextualized by the surrounding text.
\end{note}
The sole case in which abstraction differs from the B-side construction of
\Cref{subsec:b-syntax-change-repr} is the annotation construct \(\asS\). Since
annotations are used here only to resolve the type of \(\noneS\), abstraction
replaces the annotated concrete form with the primitive abstract constructor:
\[
  \aterm{\asS\ t\ \tau}{\Delta} \triangleq
  \begin{cases}
    \noneS_{\tau} & \text{if } t = \noneS, \\
    \aterm{t}{\Delta} & \text{otherwise.}
  \end{cases}
\]